\documentclass[10pt,a4paper]{amsart}
\usepackage[margin=19mm]{geometry}
\usepackage{amsmath,amssymb,mathtools}
\usepackage[scr=rsfs,cal=euler]{mathalfa}
\usepackage{xfrac}
\usepackage[unicode,hidelinks,bookmarks=false]{hyperref}
\newcommand{\ie}{i.e.\ }
\newcommand{\cf}{cf.\ }
\newcommand{\resp}{respectively\ }
\newcommand{\Subsection}{\subsection}
\providecommand{\nobreakdash}{\nobreak\hbox{-}}
\setlength{\emergencystretch}{2em}
\multlinegap0pt
\usepackage{bm}
\usepackage{bbm}
\usepackage{dsfont}
\usepackage{xspace}
\usepackage{cancel}
\usepackage{mathtools}
\usepackage{amsthm}
\makeatletter
\@ifundefined{thm}{\newtheorem{thm}{Theorem}}{}
\makeatother
\makeatletter
\expandafter\ifx\csname Thm\endcsname\relax
\newenvironment{Thm}[1][]{\refstepcounter{alphabet}%
\bigskip%
\noindent%
{\bf Theorem \Alph{alphabet}}%
\ifthenelse{\equal{#1}{}}{}{ (#1)}%
{\bf .} \itshape}{\vskip 8pt}
\fi
\expandafter\ifx\csname customthm\endcsname\relax
\newenvironment{customthm}[1]
  {\renewcommand\theinnercustomthm{#1}\innercustomthm}
  {\endinnercustomthm}
\fi
\expandafter\ifx\csname nonsec\endcsname\relax
\newenvironment{nonsec}{\bf
\setcounter{own}{\value{equation}}\addtocounter{equation}{1}
\refstepcounter{own}
\bigskip

\indent \theown.$\ \,$}{$\,$.\ \ \ }
\fi
\expandafter\ifx\csname pf\endcsname\relax
\newenvironment{pf}[1][]{%
 \vskip 3mm
 \noindent
 \ifthenelse{\equal{#1}{}}%
  {{\slshape Proof. }}%
  {{\slshape #1.} }%
 }%
{\qed\bigskip}
\fi
\makeatother
\title[Normality criterion for logharmonic mapping]{Normality criterion for logharmonic mapping}
\author{Molla Basir Ahamed, Sanju Mandal}
\date{Source: arXiv:2505.17062v1 (CC BY 4.0)}
\begin{document}
\maketitle

\noindent\emph{Source and licence.} Normality criterion for logharmonic mapping. Version arXiv:2505.17062v1 (CC BY 4.0), distributed under the Creative Commons Attribution 4.0 International licence, as recorded in the arXiv licence metadata for this exact version and shown on the abstract page https://arxiv.org/abs/2505.17062v1. Full rights notice: licenses/arxiv-2505.17062-CC-BY-4.0.txt.\par\smallskip

\noindent\textit{Editorial scope.} Counted unit: the Thm whose source label is \texttt{th-2.2} in the article (source0.tex lines 395--401, source label \texttt{th-2.2}), delivered together with the article's own complete original proof of it (source0.tex lines 402--489, one \texttt{proof} environment of 88 lines). No mathematics is written, repaired or re-derived: statement and proof are the article's own text line for line. Required context retained uncounted: th-2.1 (lines 376--378). Every definition, display and label of those lines is retained unchanged. Omitted from this package and not redistributed: 1--375, 379--394, 490--629. Nothing inside a retained excerpt is omitted or altered: every retained body is delivered exactly as the article's lines are written, so no omission receipt of any kind accompanies this package. Citations: the retained excerpts carry no citation command at all, so no bibliography entry of the article is transcribed and no reference is redistributed. Reference adaptation: none was needed and none was made; every reference inside the delivered unit resolves inside it, so no unresolved reference or citation is produced. Printed slips kept literal: no printed word, symbol, hypothesis or number of the counted statement or of its proof was altered. Layout and class: the article's own class is \texttt{amsart} with packages including amssymb, ifthen, graphicx, cite, amsmath, fontenc, inputenc, color, babel, fancyhdr, fancybox, tikz; the wrapper is the lane's pinned \texttt{amsart} editorial wrapper at 10pt on A4, which restates the theorem-like environments the retained text needs and adds the article's own macro definitions that the retained text needs (none), with extra wrapper packages (bm, bbm, dsfont, xspace, cancel, mathtools). The delivered numbering is the unit's own. Assets: no retained span contains a figure environment, a graphics-inclusion command, a PDF-page inclusion, a table or a TikZ picture; no image file is copied into this package, the package declares no asset dependency and no image file is redistributed with it. Licence: version arXiv:2505.17062v1 of this article is distributed under the Creative Commons Attribution 4.0 International licence, as recorded in the arXiv licence metadata for this exact version and shown on the abstract page https://arxiv.org/abs/2505.17062v1. A full rights notice is delivered at licenses/arxiv-2505.17062-CC-BY-4.0.txt. Characters: every retained excerpt is pure ASCII, so the delivered engine, XeLaTeX with its default Latin Modern text font, reports no missing character.
\begin{thm}\label{th-2.1}
A class $\mathcal{F}$ of logharmonic mapping $f(z)=z|z|^{2\beta}h(z) \overline{g(z)}$ in the unit disk $\mathbb{D}$ is normal if $\{f^{\#}(z):f\in\mathcal{F}\}$ uniformly locally bounded. (where $f^{\#}$ is defined in Definition 1.1)
\end{thm}

\begin{thm}\label{th-2.2}
A non-constant function $f(z)=z|z|^{2\beta}h(z) \overline{g(z)}$ logharmonic mapping in $\mathbb{D}$ is normal if and only if for each $\alpha\in(-1,\infty)$, there do not exist sequences of points $\{z_n\}\subset\mathbb{D}$ and of real numbers $\{\rho_n\}$ with $\rho_n>0$ and $\rho_n\rightarrow 0$ as $n\rightarrow\infty$ such that the functions 
\begin{align*}
	F_n(\xi):=\rho^{\alpha}_n f(z_n+\rho_n\xi)
\end{align*}
converges locally uniformly in $\mathbb{C}$ to a non-constant logharmonic mapping $F(\xi)$ satisfying $F^{\#}(\xi)\leq F^{\#}(0)= 1$.
\end{thm}

\begin{proof}
Suppose that $f$ is not normal. Then there exists a sequence $\{z^*_n\}\subset\mathbb{D}$ with $z^*_n\rightarrow 1$ as $n\rightarrow\infty$ such that 
\begin{align*}
	\left(1-|z^*_n|^2\right)\frac{|f_z(z^*_n)|+|f_{\overline{z}}(z^*_n)|}{1+|f(z^*_n)|^2}\rightarrow\infty \;\;\;\; \mbox{as}\;\;\;\;n\rightarrow\infty.
\end{align*}
Let $\{r_n\}$ be a sequence such that $|z^*_n|<r_n<1$ and
\begin{align}\label{eq-2.1}
	\left(1-\frac{|z^*_n|^2}{r^2_n}\right)\frac{|f_z(z^*_n)|+|f_{\overline{z}}(z^*_n)|}{1+|f(z^*_n)|^2}\rightarrow\infty \;\;\;\; \mbox{as}\;\;\;\;n\rightarrow\infty.
\end{align}
Now, we defined
\begin{align*}
	H_n(t,z):=\frac{\left\{\left(1-\frac{|z^*_n|^2}{r^2_n}\right)t\right\}^{1+\alpha}\left(|f_z(z)|+|f_{\overline{z}}(z)|\right)}{1 + \left\{\left(1-\frac{|z^*_n|^2}{r^2_n}\right)t\right\}^{2\alpha} |f(z)|^2},
\end{align*}
where $0<t\leq 1$ and $|z|<r_n$. Clearly, the functions $H_n(t,z)$ are continuous in $(0,1]\times\{|z|<r_n\}$. As $\alpha>-1$ and $f$ is logharmonic, it follows that
\begin{align}\label{eq-2.2}
	\lim_{t\rightarrow 0} H_n(t,z)=0.
\end{align}
On the other hand, it is easy to verify that
\begin{align}\label{eq-2.3}
	H_n(t,z)\geq t^{1+|\alpha|}\left(1-\frac{|z|^2}{r^2_n} \right)^{1+|\alpha|}  f^{\#}(z).
\end{align}
By \eqref{eq-2.1} and \eqref{eq-2.3}, we see that
\begin{align}\label{eq-2.4}
	H_n(1,z^*_n)\geq \left(1-\frac{|z^*_n|^2}{r^2_n} \right)^{1 +|\alpha|} f^{\#}(z)\rightarrow\infty \;\;\;\;\mbox{as} \;\;\;\;n\rightarrow \infty.
\end{align}
Therefore, from \eqref{eq-2.1}, \eqref{eq-2.2} and \eqref{eq-2.3}, we deduce that
\begin{align*}
	\sup_{|z|<r_n} H_n(1,z)\geq H_n(1,z^*_n)>1,\;\;\;\;\mbox{for sufficiently large} \;\; n.
\end{align*}
and
\begin{align*}
	\sup_{|z|<r_n} H_n(t,z)< 1,\;\;\;\;\mbox{for sufficiently small} \;\; t.
\end{align*}
Thus, there exists $t_n$ and $z_n$ with $0<t_n\leq 1$ and $|z_n|<r_n$ such that
\begin{align}\label{eq-2.5}
	\sup_{|z|<r_n} H_n(t_n,z)= H_n(t_n,z_n)=1.
\end{align}
Combining this with \eqref{eq-2.3} and \eqref{eq-2.5}, we obtain
\begin{align*}
	1=H_n(t_n,z_n)\geq H_n(t_n,z^*_n)\geq t^{1+|\alpha|}_n\left(1- \frac{|z^*_n|^2}{r^2_n}\right)^{1+|\alpha|}  f^{\#}(z^*_n)> t^{1+|\alpha|}_n H_n(1,z^*_n).
\end{align*}
Since the second factor $H_n(1,z^*_n)$ on the right-hand side tends to infinity by \eqref{eq-2.4}, it follows that $\lim_{n\rightarrow\infty} t_n=0$.\vspace{2mm}

Now set $\rho_n:=\left(1- \frac{|z_n|^2}{r^2_n}\right)t_n$, so that
\begin{align}\label{eq-2.6}
	\frac{\rho_n}{1-|z_n|}\leq \frac{\rho_n}{r_n -|z_n|}= \frac{(r^2_n -|z_n|^2)}{r^2_n(r_n -|z_n|)}t_n= \frac{(r_n +|z_n|)}{r^2_n} t_n \rightarrow 0 \;\;\;\;\mbox{as} \;\;\;\;n\rightarrow \infty.
\end{align}
Then the functions $F_n(\xi)=\rho^{\alpha}_n f(z_n+\rho_n\xi)$ is defined for $|\xi|<R_n=\frac{1-|z_n|}{\rho_n}$. Moreover, a simple computation yields
\begin{align}\label{eq-2.7}
	F^{\#}_n(\xi)= \frac{\rho^{1+\alpha}_n\left(|f_z(z_n+\rho_n\xi)| + |f_{\overline{z}}(z_n+\rho_n\xi)|\right)}{1+\rho^{2\alpha}_n |f(z_n+\rho_n\xi)|^2}.
\end{align}
In view of \eqref{eq-2.5} and \eqref{eq-2.7}, we obtain
\begin{align}\label{eq-2.8}
	F^{\#}_n(0) \nonumber&= \frac{\rho^{1+\alpha}_n\left(|f_z(z_n)| + |f_{\overline{z}}(z_n)|\right)}{1+\rho^{2\alpha}_n |f(z_n)|^2}\\&= \frac{\left\{\left(1- \frac{|z_n|^2}{r^2_n}\right)t_n \right\} ^{1+\alpha}\left(|f_z(z_n)| + |f_{\overline{z}}(z_n)|\right)} {1+\left\{\left(1- \frac{|z_n|^2}{r^2_n}\right)t_n \right\}^{2\alpha} |f(z_n)|^2} = H_n(t_n,z_n)=1.
\end{align}
It follows from \eqref{eq-2.6} that $\left(1- \frac{|z_n|^2}{r^2_n} \right)/\left(1- \frac{|z_n+ \rho_n\xi|^2}{r^2_n}\right)$ tends uniformly to $1$ as $n\rightarrow\infty$ on compact subsets of $\mathbb{C}$. Let us choose $\epsilon_n>0$ with $\epsilon_n\rightarrow 0$ as $n\rightarrow\infty$ such that
\begin{align}\label{eq-2.9}
	(1-\epsilon_n)\left(1- \frac{|z_n+ \rho_n\xi|^2}{r^2_n}\right)t_n \leq\rho_n\leq (1+\epsilon_n)\left(1- \frac{|z_n+ \rho_n\xi|^2} {r^2_n}\right)t_n.
\end{align}
Therefore, from \eqref{eq-2.5}, \eqref{eq-2.7} and \eqref{eq-2.9}, we obtain
\begin{align}\label{eq-2.10}
	F^{\#}_n(\xi)\nonumber&\leq \frac{(1+\epsilon_n)^{1+\alpha} \left\{\left(1- \frac{|z_n+ \rho_n\xi|^2} {r^2_n}\right)t_n \right\}^{1+\alpha} \left(|f_z(z_n+\rho_n\xi)| + |f_{\overline{z}} (z_n+\rho_n\xi)| \right)}{1+(1-\epsilon_n) ^{2\alpha} \left\{\left(1- \frac{|z_n+ \rho_n\xi|^2}{r^2_n}\right) t_n\right\}^{1+\alpha} |f(z_n+ \rho_n\xi)|^2} \\&\nonumber\leq \frac{(1+\epsilon_n)^{1+\alpha}\left\{\left(1- \frac{|z_n+ \rho_n\xi|^2} {r^2_n}\right)t_n\right\}^{1+\alpha} \left(|f_z(z_n +\rho_n\xi)| + |f_{\overline{z}}(z_n+\rho_n\xi)| \right)}{(1- \epsilon_n)^{2\alpha} \left\{\left(1- \frac{|z_n+ \rho_n\xi|^2} {r^2_n}\right)t_n\right\}^{1+\alpha} |f(z_n+ \rho_n\xi)|^2} \\&\nonumber=\frac{(1+\epsilon_n)^{1+\alpha}}{(1-\epsilon_n)^{2\alpha}} \cdot H_n(t_n, z_n+\rho_n\xi) \cdot \frac{1+ \left\{\left(1- \frac{|z_n+ \rho_n\xi|^2}{r^2_n}\right)t_n\right\}^{1+\alpha} |f(z_n+ \rho_n\xi)|^2}{\left\{\left(1- \frac{|z_n+ \rho_n\xi|^2} {r^2_n}\right)t_n\right\}^{1+\alpha} |f(z_n+ \rho_n\xi)|^2} \\&\leq\frac{(1+\epsilon_n)^{1+\alpha}} {(1-\epsilon_n)^{2\alpha}} \cdot H_n(t_n, z_n+\rho_n\xi) \rightarrow 1 \;\;\;\;\mbox{as} \;\;\;\;n\rightarrow \infty.
\end{align}
Thus, by Marty's Theorem \ref{th-2.1}, we conclude that the sequence $\{F_n(\xi)\}$ is normal. We may assume that $\{F_n(\xi)\}$ converges locally uniformly in $\mathbb{C}$. Then, the limit function $\{F(\xi)\}$ is logharmonic in $\mathbb{C}$. In view of \eqref{eq-2.8}
and \eqref{eq-2.10} it follows that $F^{\#}(\xi)\leq F^{\#}(0)= 1\neq 0$, which implies that $F$ is a non-constant logharmonic mapping. \vspace{2mm}

Next, we prove the necessary part of the theorem. Let $f$ be normal in $\mathbb{D}$. Again, we recall that the functions $F_n(\xi)$ given by
\begin{align*}
	F_n(\xi)=\rho^{\alpha}_n f(z_n+\rho_n\xi)
\end{align*}
are defined for $|\xi|<\frac{1-|z_n|}{\rho_n}$. Again, we observe that
\begin{align*}
	\frac{\rho^{1+\alpha}_n}{1-|z_n|}\leq \frac{\rho^{1+\alpha}_n}{r_n -|z_n|}&= \frac{(r^2_n -|z_n|^2)^{1+\alpha}}{r^{2(1+\alpha)}_n(r_n -|z_n|)}t^{1+\alpha}_n \\&= \frac{(r_n +|z_n|)^{1+\alpha}(r_n -|z_n|)^{\alpha}}{r^{2(1+\alpha)}_n} t^{1+\alpha}_n \\&< \frac{r^{\alpha}_n(r_n +|z_n|)^{1+\alpha}}{r^{2(1+\alpha)}_n} t^{1+\alpha}_n \\&= \frac{(r_n +|z_n|)^{1+\alpha}}{r^{\alpha +2}_n} t^{1+\alpha}_n \\&< \frac{2^{1+\alpha}}{r_n} t^{1+\alpha}_n \rightarrow 0 \;\;\;\;\mbox{as} \;\;\;\;n\rightarrow \infty,
\end{align*}
which implies that
\begin{align*}
	\frac{\rho^{1+\alpha}_n}{1-|z_n|-\rho_n |\xi|}\rightarrow 0 \;\;\;\;\mbox{as} \;\;\;\;n\rightarrow \infty,\;\;\mbox{for}\;\; |\xi|<\frac{1-|z_n|}{\rho_n}.
\end{align*}
Since
\begin{align*}
	F^{\#}_n(\xi)&= \frac{\rho^{1+\alpha}_n\left(|f_z(z_n+\rho_n\xi)| + |f_{\overline{z}}(z_n+\rho_n\xi)|\right)}{1+\rho^{2\alpha}_n |f(z_n+\rho_n\xi)|^2} \\&\leq \frac{\rho^{1+\alpha}_n \left(|f_z(z_n+\rho_n\xi)| + |f_{\overline{z}}(z_n+ \rho_n\xi)| \right)}{1+ |f(z_n+\rho_n\xi)|^2} \\&= \rho^{1+ \alpha}_n f^{\#}(z_n+\rho_n\xi)\\&\leq \frac{\rho^{1+\alpha}_n} {1-|z_n|-\rho_n |\xi|} \left(1 -|z_n+\rho_n\xi|^2\right) f^{\#}(z_n+\rho_n\xi)
\end{align*}
and $f$ is normal, we have
\begin{align*}
	\left(1 -|z_n+\rho_n\xi|^2\right) f^{\#}(z_n+\rho_n\xi)<\infty.
\end{align*}
Therefore, it is easy to see that $F^{\#}_n(\xi)\rightarrow 0$ as $n\rightarrow\infty$ and thus, $F^{\#}(\xi)=0$ for all $\xi\in \mathbb{C}$, so that $F(\xi)$ is a constant, which is a contradiction. This completes the proof.
\end{proof}

\end{document}
